\chapter{备选与观察：为什么暂不升级}

\section{中国船舶：质量最好的零权备选}

中国船舶2025年收入1,519.78亿元、归母78.48亿元；2026Q1收入433.12亿元、归母48.32亿元、扣非47.67亿元、经营现金流81.52亿元。H1预告归母92--110亿元、扣非90--108亿元。利润、扣非与现金三者同时改善，是六只深度标的中最扎实的报表组合。

原始券商报告预计2026--2028年归母净利润181.98/235.38/299.43亿元，EPS 2.42/3.13/3.98元，并判断交付排期延伸至2028--2030年。订单能见度和高价订单进入交付构成核心逻辑。需要注意的是，2025年资产整合和同一控制下合并会扭曲同比增速，不能直接把并表增长套入高成长PE。

现价33.02元约对应2027E EPS 10.55倍，翻倍需要21.10倍。AStock牛市情景68.9元可以超过翻倍门槛，但当前收集的五份原始报告均没有显式目标价。不能从报告列示的当前PE反推券商目标，因此该情景保持零外部权重。

\begin{keyinsight}[升级条件]
若三季报毛利率维持约16\%以上、经营现金流持续为正、高价与中高端船交付继续，同时取得当前、显式、可审计的外部目标，中国船舶可以重新竞争条件核心。它的优势是质量，缺点是五个月内将倍数翻倍的难度。
\end{keyinsight}

\section{盛新锂能：商品周期尾部}

盛新锂能2025年归母亏损8.88亿元，2026Q1归母转为4.64亿元，H1预告归母10--12亿元、扣非13--15亿元。反转由锂盐价格、出货、自有矿与印尼工厂共同驱动，券商预测2026E EPS 2.45元、2027E EPS 3.28元。

现价27.65元翻倍至55.30元，需要22.57倍2026E EPS；若用2027E EPS，仍需16.86倍。AStock牛市情景59.0元能够翻倍，但原始报告没有显式目标，且2026Q1经营现金流为-6.67亿元。量价与自供比例必须和现金共同验证，单一锂价上涨不足以升级。

需要观察Q3出货是否接近3万吨、综合吨利是否接近2万元、印尼利用率是否继续提升、经营现金流是否转正。若全年出货接近12万吨、自有矿贡献和现金均可核验，才有资格从周期备选上调。

\section{江波龙：利润很强，但分母可能在峰值}

江波龙2026Q1收入99.09亿元、归母38.62亿元、扣非39.43亿元，毛利率55.53\%；H1预告收入220--250亿元、归母92--110亿元。利润增长极强，但Q1经营现金流为-28.75亿元，FY2025经营现金流也为-12.01亿元。

国信证券预测2026E净利润110.97亿元、EPS 26.47元，而2027E净利润仅增长约1.9\%。另一份原始报告对2026E净利润预测高达199.69亿元，两份报告在类似收入预测下出现88.72亿元利润差，说明晶圆价格、库存成本层与毛利假设主导模型。如此大的分歧不适合建立精确目标。

股价从7月9日620元跌至7月22日388.45元，十个交易日下跌37.35\%。现价翻倍需要约29.35倍2026E EPS。若2027利润不增长，市场必须在峰值分母上给出接近30倍PE，难度与回撤风险都很高。

\begin{riskbox}[核心阻断]
公司尚未披露bit出货、产品ASP、库存收益拆分、封测利用率、下游订单转量与客户收入。平台认证与LTA/MOU能够说明技术和供给关系，却不能证明可持续单位经济。只有单季经营现金流转正、库存周转改善、H2毛利稳定并出现2027年两位数利润增长，才应重建估值。
\end{riskbox}

\section{天华新能：盈利强，但当前价格已使翻倍门槛更高}

天华新能H1预告归母22--24亿元、扣非21.62--23.62亿元，Q1毛利率46.99\%，盈利恢复显著。券商预测2026E EPS约5.70元，AStock情景为56.4/79.8/102.4元。以现价55.88元计算，牛市空间约83.3\%，仍未达到翻倍所需的111.76元。

公司与宁德时代的股权关系可说明产业协同，不能自动推导采购订单和价格溢价；未投产矿山与在建二期产能也不进入2026现金流。现阶段它是有效的锂周期观察标的，但不是本题要求的翻倍候选。

\begin{exhibitbox}[Exhibit 7：零权备选的阻断项]
\small
\begin{tabularx}{\exhibitboxwidth}{L{2.4cm} R{1.5cm} R{1.5cm} X L{3.2cm}}
\toprule
标的 & 牛市空间 & Q1 OCF & 主要优势 & 发布阻断 \\
\midrule
中国船舶 & +108.7\% & +81.52亿元 & 订单、交付、现金质量 & 无显式外部目标；并表归一化 \\
盛新锂能 & +113.5\% & -6.67亿元 & 锂价、自有矿、印尼增量 & 无显式目标；商品与现金风险 \\
江波龙 & +100.8\% & -28.75亿元 & 存储价格与产品升级 & 峰值分母、预测分歧、现金缺口 \\
天华新能 & +83.3\% & -1.83亿元 & 锂盐量价与资源布局 & 牛市区间本身不翻倍 \\
\bottomrule
\end{tabularx}
\sourcenote{公司报告、H1预告、原始券商报告、AStock情景模型。}
\end{exhibitbox}

\clearpage
